\chapter{Introduction}\label{ch:introduction}
% Content cycle: C2. Intent: state the portability problem; motivate the
% algorithm/schedule split; give the thesis statement + contributions; bound
% the scope (the algorithm class); lay out the dissertation structure.

\section{The portability problem}\label{sec:intro-problem}
% Porting an algorithm across CPU threads / MPI cluster / embedded MCU means
% rewriting it: decomposition + IO semantics are hand-written, per platform.

\section{Motivation}\label{sec:intro-motivation}
% IO is a second-class citizen (compute scheduling has Halide/TVM/Exo; IO
% does not); hardware churns faster than algorithms; decomposition is the
% expensive, error-prone part.

\section{Thesis statement and contributions}\label{sec:intro-contributions}
% Algorithm/schedule separation; IO as a first-class schedule directive; one
% Petri-net IR with a compile-time soundness gate; falsifiable genericity via
% a cross-backend bit-identical differential test.

\section{Scope and non-goals}\label{sec:intro-scope}
% The affine/static/single-assignment algorithm class; what is deliberately
% out (data-dependent writes, recursion, dynamic scheduling).

\section{Structure of this dissertation}\label{sec:intro-structure}

\todo{Write Chapter 1 (cycle C2).}
